% GENERATED FILE. Edit canonical lessons, then run npm run generate:books.
% canonical-source-hash: fnv1a64:b4daa29ac0124372
% canonical-lessons: KA-C140-gadda, KA-C140-pada, KA-C140-udda, KA-C140-gidda, KA-C140-bhara

\chapter{Beard, Foot, Long, Short, Heavy}
\label{ch:ka-beard-foot-long-short-heavy}
\hlchaptermodality{140}

\begin{chapteropening}
\textbf{By the end of this chapter:} \emph{I can say a beard, a foot, long, short and heavy.}

Complete the last lesson of chapter 140: i can say a beard, a foot, long, short and heavy.
\end{chapteropening}

\section[gaḍḍa]{\kn{ಗಡ್ಡ} (gaḍḍa) — a beard}

\label{lesson:KA-C140-gadda}



\begin{quote}
Before the new one: say the Kannada for tears, then the Kannada for a moustache.
\end{quote}

\subsection*{You'll want to know: \kn{ಗಡ್ಡ}}
\textbf{\kn{ಗಡ್ಡ}} — \emph{gaḍḍa} — \textquotedblleft{}a beard\textquotedblright{}.

\begin{grammarlens}[title={What you've built}]
One more word for this chapter.
\end{grammarlens}

\subsection*{Your turn}
\begin{itemize}
\raggedright
  \item \emph{Say it:} \emph{gaḍḍa}
  \item \emph{Say it:} \emph{gaḍḍa}, once more
  \item \emph{Say it:} say \emph{gaḍḍa}
  \item \emph{Recall:} say the Kannada for tears, then the Kannada for a moustache, then say \emph{gaḍḍa} again
\end{itemize}

\begin{tcolorbox}[breakable,colback=teal!4,colframe=teal!35!black,title={Before you move on}]
What does \kn{ಗಡ್ಡ} mean? (\textquotedblleft{}A beard\textquotedblright{}.) Say it once more.
\end{tcolorbox}
\section[pāda]{\kn{ಪಾದ} (pāda) — a foot}

\label{lesson:KA-C140-pada}



\begin{quote}
Before the new one: say the Kannada for a moustache, then the Kannada for a beard.
\end{quote}

\subsection*{You'll want to know: \kn{ಪಾದ}}
\textbf{\kn{ಪಾದ}} — \emph{pāda} — \textquotedblleft{}a foot\textquotedblright{}.

\begin{grammarlens}[title={What you've built}]
One more word for this chapter.
\end{grammarlens}

\subsection*{Your turn}
\begin{itemize}
\raggedright
  \item \emph{Say it:} \emph{pāda}
  \item \emph{Say it:} \emph{pāda}, once more
  \item \emph{Say it:} say \emph{pāda}
  \item \emph{Recall:} say the Kannada for a moustache, then the Kannada for a beard, then say \emph{pāda} again
\end{itemize}

\begin{tcolorbox}[breakable,colback=teal!4,colframe=teal!35!black,title={Before you move on}]
What does \kn{ಪಾದ} mean? (\textquotedblleft{}A foot\textquotedblright{}.) Say it once more.
\end{tcolorbox}
\section[udda]{\kn{ಉದ್ದ} (udda) — long, tall}

\label{lesson:KA-C140-udda}



\begin{quote}
Before the new one: say the Kannada for a beard, then the Kannada for a foot.

An adjective goes before the noun and does not change: \textbf{\kn{ದೊಡ್ಡ} \kn{ಮರ}}, \textbf{\kn{ಚಿಕ್ಕ} \kn{ಮರ}}.
\end{quote}

\subsection*{You'll want to know: \kn{ಉದ್ದ}}
\textbf{\kn{ಉದ್ದ}} — \emph{udda} — \textquotedblleft{}long, tall\textquotedblright{}.

\begin{grammarlens}[title={What you've built}]
One more word for this chapter.
\end{grammarlens}

\subsection*{Your turn}
\begin{itemize}
\raggedright
  \item \emph{Say it:} \emph{udda}
  \item \emph{Say it:} \emph{udda}, once more
  \item \emph{Say it:} say \emph{udda}
  \item \emph{Recall:} say the Kannada for a beard, then the Kannada for a foot, then say \emph{udda} again
\end{itemize}

\begin{tcolorbox}[breakable,colback=teal!4,colframe=teal!35!black,title={Before you move on}]
What does \kn{ಉದ್ದ} mean? (\textquotedblleft{}Long, tall\textquotedblright{}.) Say it once more.
\end{tcolorbox}
\section[giḍḍa]{\kn{ಗಿಡ್ಡ} (giḍḍa) — short}

\label{lesson:KA-C140-gidda}



\begin{quote}
Before the new one: say the Kannada for a foot, then the Kannada for long.
\end{quote}

\subsection*{You'll want to know: \kn{ಗಿಡ್ಡ}}
\textbf{\kn{ಗಿಡ್ಡ}} — \emph{giḍḍa} — \textquotedblleft{}short\textquotedblright{}.

\begin{grammarlens}[title={What you've built}]
One more word for this chapter.
\end{grammarlens}

\subsection*{Your turn}
\begin{itemize}
\raggedright
  \item \emph{Say it:} \emph{giḍḍa}
  \item \emph{Say it:} \emph{giḍḍa}, once more
  \item \emph{Say it:} say \emph{giḍḍa}
  \item \emph{Recall:} say the Kannada for a foot, then the Kannada for long, then say \emph{giḍḍa} again
\end{itemize}

\begin{tcolorbox}[breakable,colback=teal!4,colframe=teal!35!black,title={Before you move on}]
What does \kn{ಗಿಡ್ಡ} mean? (\textquotedblleft{}Short\textquotedblright{}.) Say it once more.
\end{tcolorbox}
\section[bhāra]{\kn{ಭಾರ} (bhāra) — heavy}

\label{lesson:KA-C140-bhara}



\begin{quote}
Before the new one: say the Kannada for long, then the Kannada for short.
\end{quote}

\subsection*{You'll want to know: \kn{ಭಾರ}}
\textbf{\kn{ಭಾರ}} — \emph{bhāra} — \textquotedblleft{}heavy\textquotedblright{}.

\begin{grammarlens}[title={What you've built}]
One more word for this chapter.
\end{grammarlens}

\subsection*{Your turn}
\begin{itemize}
\raggedright
  \item \emph{Say it:} \emph{bhāra}
  \item \emph{Say it:} \emph{bhāra}, once more
  \item \emph{Say it:} say \emph{bhāra}
  \item \emph{Recall:} say the Kannada for long, then the Kannada for short, then say \emph{bhāra} again
\end{itemize}

\begin{tcolorbox}[breakable,colback=teal!4,colframe=teal!35!black,title={Before you move on}]
What does \kn{ಭಾರ} mean? (\textquotedblleft{}Heavy\textquotedblright{}.) Say it once more.
\end{tcolorbox}
